\subsection{扬州的檄文与运河上的兵粮}
684年秋，徐敬业在扬州举兵时，运河并不是一条为首都服务的直线，而是一串需要船户、仓吏、渡口和沿岸州县维持的节点。\eventanchor{ST-684-LI-JINGYE}{徐敬业扬州起兵（684）}把反武后的主张写成檄文，正说明反对者必须借用唐的官爵、宗室与公共文辞来召集同盟；檄文能告诉我们政治语言如何安排敌我，不能替我们统计有多少乡里因它投入战斗。朝廷的反应同样有两层：洛阳可以任命将帅、宣布罪名，却仍须让粮、船、武器和命令顺水陆抵达江都附近。胜败并非仅由一篇文字或一位统帅决定。扬州附近的水网让集结与移动成为可能，也使一支军队离开补给点后容易被切断。两《唐书》和《资治通鉴》保留了起兵、讨平的年序；它们更擅长记录中央如何为胜利命名，较难留下被临时征船、守渡或避兵者的连续经历。

这场失败没有使江淮从政治中退场。恰相反，朝廷此后更清楚地看到，国号的合法性要经过地方交通才能变成可执行的命令。把徐敬业的行动写成孤臣的浪漫反抗，会遮蔽他利用的军政遗产；把它写成单纯的地方叛乱，又会遮蔽其针对继承与临朝的公开论证。事件的制度得失在于：中央暂时取得了惩办与任命的主动，却也把控制运输节点、审查地方人事和取得可靠军队变成长期成本。后来武周对河北、河西与江淮的不同处置，都在回答同一个问题：一纸诏书抵达之后，谁愿意替它保管仓门、修补道路并承担风险。

\subsection{洛阳、受命与寺院的可见政治}
天授元年改国号的仪式在洛阳完成，\eventref{ST-690-ZHOU-FOUNDATION}{周的建立}因而首先是一场关于城、宫门、礼仪与文书的行动。官员的请愿、改换年号、宗室的处置和新官名并置在同一过程里；它们不是装点权力的背景，而是让命令被抄写、传递、存档和认可的手段。史书所说的群臣拥戴带有新朝与复唐叙事反复改写的痕迹，不能把它当作所有官员和州县的内心投票。不过，改号、称帝和随后的人事调整本身足以说明唐的旧有行政网络并未消失；武周必须使用它，又必须重新安排其忠诚和上升通道。

691年围绕《大云经》的政治使用，\eventanchor{ST-691-GREAT-CLOUD-SUTRA}{《大云经》与新朝合法性（691）}把佛教经义带入受命语言。经文、注疏和僧侣网络能够提供一种解释新朝的词汇，却不等于所有寺院被朝廷直接操控，更不能证明某种预言为当时社会普遍相信。僧传和后出的正史尤其倾向把宗教人物写成护法者或谄谀者；造像题记、写经和寺院遗存只能照亮有限地点的供养关系。可被较稳妥地观察到的，是国家礼仪与宗教组织在洛阳及若干地方相遇：授予、翻译、题记、布施与官府保护，使权威既能借宗教显形，也受到寺院财产、地方施主和跨区流动的限制。

因此，不能把武周的佛教政策缩成“利用宗教”的一句判断。对朝廷而言，符瑞语言提供了处理性别、宗统与天命断裂的资源；对僧侣、译者和施主而言，国家的承认可能带来声望、通行和财产保护，也可能把他们卷入竞争。宗教在这里不是国家之外的温和领域，而是一处能聚集纸张、知识、捐施与人际关系的制度空间。它的代价是：当宫廷合法性需要借它说话时，宗教组织也更容易成为后来的政治批评和清算对象。

\subsection{四镇、北庭与不能画直的边界}
692年恢复安西四镇，\\eventref{ST-692-ANXI-RECOVERY}{恢复安西四镇（692）}，在两《唐书》的叙事中常呈为唐军重返西域的结果。若把这句话当作一幅完整疆界图，便会漏掉龟兹、于阗、疏勒、碎叶之间的距离、绿洲政治和不同的驻军条件。夺回或重新控制据点与维持据点是两件事：前者需要一次军事推进，后者要有粮秣、马匹、翻译、文书、地方中介和不断接替的守军。汉文传记记录的城镇顺序、军功和官名是重要线索，却不足以复原每个绿洲的选择。吐鲁番文书和考古所见的行政痕迹能使某些物资、人名和程序可见，但其地区性正告诫我们不可将一个绿洲的实践投射成“唐制”的全貌。

702年置北庭都护府，\eventanchor{ST-702-BEITING}{北庭都护府的设置（702）}，把庭州与天山北路推到更显眼的位置。这个动作不是把草原和山地纳入一张已完成的行政地图，而是为沿路军政、使节与贸易设定一个可以交涉的枢纽。当地势、季节和草场限制行军，中央的册命则限制什么人能以何种身份领取军资、传递消息或同邻近政权谈判。都护府的名字是确定的，控制的深浅却随道路安全、联盟和补给而变化。边疆史最容易在“失地”“收复”的词语中抹去这种变化；较好的读法是跟着具体城镇、关隘和中转点看国家能力怎样被暂时接上，又在哪里断裂。

武周在西北的选择也反过来影响东部。投入河西、安西与北庭的资源不能同时解决所有河北军情；相反，东北的起事会迫使朝廷重新分配将领和兵粮。这不是说某次西域经营必然造成契丹冲突，而是说同一国家的道路与财政从来不是无限的。它的制度得失在于建立了可用的外部接口，代价则是不断把内地税役、地方运输和边地中介纳入更长的供给链。

\subsection{营州的契丹起事与河北的道路}
万岁通天元年，李尽忠、孙万荣在营州起事，\\eventref{ST-695-KHITAN-REBELLION}{营州契丹起事（695）}。正史以周廷的军报来组织这一段历史，因此“叛”“降”“平”首先是朝廷分类，而不是对契丹部众目标的透明记录。起事和周军初期失利可以确定；首领之间的权力关系、部众构成以及每一座城镇的态度，却不能由胜者的记录补全。营州周边并非单一的边缘地带：它连接辽西道路、州县军镇、草原移动和东北各群体的贸易、婚姻与冲突。消息、粮草和援军要穿过这些联系，中央才有可能把洛阳的决定转成战场行动。

697年孙万荣败死，起事告一段落，\eventanchor{ST-697-KHITAN-SUPPRESSION}{契丹战争的结束（697）}；这也不是河北从此恢复静止的时刻。后突厥势力的扩展、边镇人事和地方征发继续改变东北的条件。战役结局使朝廷能重新表彰将领、布置守备，不能量出沿线家庭损失了多少劳力，也不能把被迫迁动者写成没有选择的背景。边地材料的沉默在这里尤其重要：汉文帝纪留下命令和军功，外部传统与后世族群叙事又各有自己的政治时间。叙事所能做的，是避免把后来的民族名称和稳定边界直接压回七世纪末的移动人群。

对武周而言，这场冲突揭示的不是“北方天生不服”，而是行政、交通和联盟的脆弱接口。为回应起事而增加的军事需求，会向河北州县、道路和仓储传递压力；中央若只通过惩罚来理解它，就会忽略中介者、地形和资源如何塑造行动。此处的中期影响，是东北军事经验与官员网络继续进入复唐和开元的用人结构，成为后来边镇集中问题的一部分，而非一条直线的宿命。

\subsection{还唐不是把十五年倒带}
698年武则天召回庐陵王李显、复立为太子，\eventanchor{ST-698-LI-XIAN-HEIR}{李显复立太子（698）}，表明继承安排已重新转向李唐宗室。这个决定并不证明宫廷矛盾已经消失；恰恰相反，它显示新朝需要在官僚、人事与宗室之间寻找可被接受的继承形式。705年张柬之等发动政变，\\eventref{ST-705-TANG-RESTORATION}{神龙政变与复唐（705）}，逼迫武则天退位，恢复唐号。政变的日期和结果可靠，参与者分工、宫禁中的行动顺序以及清算范围在纪传中却有异文。读者不必从这些差异中挑选一个戏剧性版本；更值得追问的是，为什么宫廷中的联盟可以在特定时刻调动禁军、控制门禁并让诏书迅速取得效力。

复唐改变了国号和中枢控制的归属，却不能使地方官、寺院、家族和边镇从武周年号下的日常行政中突然抽离。墓志的年号、官衔和迁葬文字有时能校对个体在转换中的位置，首先仍是家族选择留下的公开记忆，不能据此测量整个地区的“民心”。复唐的实际速度取决于谁传递诏命、谁继续管理仓库与户籍、谁能在地方调动武装。这种连续并不减损政变的重要性，反而说明政权更替为何总要依赖旧有行政劳动。

武周留下的制度问题也没有随着国号消散：宫廷怎样处置宗室，边疆怎样筹饷，宗教组织同国家礼仪如何相处，地方人事怎样既可被提拔又可被监控。复唐把合法性的语汇换回李氏，却没有解除这些压力。随后数次宫廷冲突之所以显得急促，正是因为继承和禁军控制仍是决定谁能把命令送出宫门的条件。

\subsection{唐隆政变的宫门与继承的成本}
707年太子李重俊举兵试图诛杀武三思而失败，\eventanchor{ST-707-LI-CHONGJUN}{李重俊起兵（707）}，显示复唐后的宗室并未获得稳定的政治位置。史书对太子、武三思和支持者的褒贬浓重，不能把其中的罪状直接当作完整动机。仍可看见的是，继承人若要改变局势，必须取得宫禁内外的兵力、官员与消息；一旦这些联系断裂，血统本身不能维持行动。太子失败带来的不是抽象的教训，而是新的恐惧：谁能接近皇帝，谁掌握禁军，谁能把私人冲突转成公开诏令。

710年六月，李隆基与太平公主发动唐隆政变，\eventanchor{ST-710-TANGLONG-COUP}{唐隆政变（710）}，诛韦后集团、拥立睿宗。关于宫禁行动、死亡方式和参与名单的细节，各史并不一致；这些差异提醒我们，后来胜者非常关心把暴力写成恢复秩序的必然步骤。不过，政变需要在极短时间内占住门禁、联络禁军、控制文书和宣布新的继承人，这些物质而制度化的条件比人物性格更能解释其可行性。宫廷政变也不是只发生在宫墙之内：一旦政令传出，州县和边镇须辨认哪个年号、哪份诏书、哪位使者可以被承认。

712年睿宗让位于李隆基，\\eventref{ST-712-XUANZONG-ACCESSION}{玄宗即位（712）}，父子之间的权力并未立即完全分开。713年，玄宗诛太平公主党羽、太平公主自尽，\\eventref{ST-713-TAIPING}{平太平公主（713）}，并存的权力中心才告终结。这里同样要警惕胜者档案：所谓谋反与清洗名单不能不经辨析地成为事实。但事件留下的制度后果清楚：玄宗把禁军、诏令和人事的最后解释权收归一处，使中枢能够更连贯地处理户籍、漕运和边防；代价则是后续行政能力会更依赖少数近臣、将领与中枢网络的平衡。

\subsection{括户的笔与村落的门}
开元九年宇文融主持括户检田，\\eventref{ST-721-YUWEN-RONG}{宇文融括户（721）}。从朝廷的角度看，漏户、逃户和隐田都是应被重新纳入名册的资源；从州县办事者的角度看，它们意味着要在旧籍、亲属关系、流动人口、豪强占有和实际田亩之间作出可执行的判断。一个名字被写入簿册，可能改变的是赋役、差役和地方交涉，而不只是中央账面增加的一项数字。两《唐书》、通鉴和传记保存了任命、奏报与争论；其新增户口和田亩是行政口径，不能直接等同于全国人口或财富。

敦煌和吐鲁番的户籍、计帐等文书让我们知道登记并非抽象词语：纸张、格式、担保、地名与人员身份都需要被具体写下。然而它们属于不同时间和地区，不能用来证明开元各地一律按同一节奏清查。制度的力量正体现在这种不整齐中：中央提出分类目标，地方社会则在迁徙、依附、婚姻、土地与季节性生计中与之周旋。括户解决的是朝廷在和平时期重新发现税源的需要；它的副作用是把过去模糊或可协商的身份推到更直接的征收关系中，也使地方官更深地介入家庭的移动和财产叙述。

\subsection{封禅、漕运与繁荣的两种舞台}
725年玄宗封禅泰山，\eventanchor{ST-725-TAISHAN}{开元封禅（725）}。仪式所展示的不是已经无需劳动的富足，而是朝廷有能力调集礼官、随从、道路和供给，将一套天下秩序放到山岳之上。礼仪中的四方服从与仓廪充实是政治语言，不能据此推断每个地区都同样安定。与666年高宗、武后封禅相比，开元的泰山也提示皇权必须周期性地通过地理、礼制与可见的组织来确认自身；它并不能替代对河北、江淮、河西或岭南具体处境的调查。

733年裴耀卿整顿江淮漕运，\\eventref{ST-733-PEI-YAOQING}{裴耀卿转运（733）}，提供了另一种更不华丽的国家舞台。关中的供给不是从地图上自动流来：水位、季节、船只、仓场、转运点和沿线劳力都会决定粮食能否按时抵达。纪传和《通典》记录了整顿的名称与朝廷评价，漕粮数额和效率却不能不经检验地化为全国经济状况。沿线人群面对的也并非单纯“国家需要”：有人参与雇役与交易，有人负担征发或等待水道通行，地方官须在仓储、税收和民生之间处理冲突。

封禅与转运相连，并不是因为一件导致另一件，而是因为两者都暴露开元秩序依赖的中介。山上的礼仪需要道路和供给；长安的粮食需要江淮的水、船和账簿。国家能力在此不是一句“盛世”，而是许多能被接续、也可能被阻断的劳动。\endinput
